%Chargement des paquets
\usepackage{amsmath}
\usepackage{amsfonts}
\usepackage{amssymb}
\usepackage{enumerate}
\usepackage{mathtools}
\usepackage{bbm}
\usepackage{xparse, etoolbox}
\usepackage{enumerate}
\usepackage{enumitem}
\usepackage{mathabx}
\usepackage{babel}[french]

%environment
\usepackage{amsthm}
\usepackage{keytheorems}
\usepackage{hyperref} 

%Interface théorème
\newkeytheorem{lem}[
	name = Lemme
]

\newkeytheorem{prop}[
	name = Proposition
]

\newkeytheorem{thm}[
	name = Théorème
]

\newkeytheorem{coro}[
	name = Corollaire
]

\renewcommand*{\proofname}{Démonstration}

\theoremstyle{definition}
\newtheorem{defn}{Définition}

\theoremstyle{definition}
\newtheorem*{nota}{Notation}

\theoremstyle{definition}
\newtheorem*{limi}{Liminaire}

\theoremstyle{remark}
\newtheorem{rmq}{Remarque}

\theoremstyle{plain}
\newtheorem*{fait}{Fait}


%Ensembles de nombres
\newcommand{\N} {\mathbb{N}}
\newcommand{\Ne}{\N^\ast}
\newcommand{\Z} {\mathbb{Z}}
\newcommand{\Q} {\mathbb{Q}}
\newcommand{\R} {\mathbb{R}}
\newcommand{\Rb}{\overline{\mathbb{R}}}
\newcommand{\Rp}{\R_+}
\newcommand{\Rm}{\R_-}
\newcommand{\K} {\mathbb{K}}
\newcommand{\Ket} {\mathbb{K^{\ast}}}
\newcommand{\Cx}{\mathbb{C}}

%Ensembles, tribus
\newcommand{\A}{\mathbb{A}}
\renewcommand{\P}{\mathcal{P}}
\newcommand{\T}{\mathcal{T}}
\newcommand{\F}{\mathcal{F}}
\newcommand{\C} {\mathcal{C}}
\newcommand{\ouv}{\mathcal{O}}
\newcommand{\B}{\mathcal{B}}
\newcommand{\M}{\mathcal{M}}
\newcommand{\etag}{\mathcal{E}}
\newcommand{\etagOp}{\etag^{+}(\Omega)}
\newcommand{\etagO}{\etag(\Omega)}
%%Lp avant quotientage
\newcommand{\Lic}{\mathcal{L}^1}
\newcommand{\Lpc}{\mathcal{L}^p}
\newcommand{\Linfc}{\mathcal{L}^{\infty}}
\newcommand{\Lifull}{\Lic(\Omega, \B, m)}
\newcommand{\Lpfull}{\Lpc(\Omega, \B, m)}
\newcommand{\Linffull}{\Linfc(\Omega, \B, m)}
%%Lp après quotientage
\newcommand{\Li}{L^1}
\newcommand{\Ld}{L^2}
\newcommand{\Lp}{L^p}
\newcommand{\Linf}{L^{\infty}}
\newcommand{\Listd}{\Li(\Omega, m)} 
\newcommand{\Ldstd}{\Ld(\Omega, m)} 
\newcommand{\Lpstd}{\Lp(\Omega, m)} 

%Quelques opérateurs
\DeclareMathOperator{\codim}{codim}
\DeclareMathOperator{\rg}{rg}
\DeclareMathOperator{\tr}{tr}
\DeclareMathOperator{\vect}{Vect}
\DeclareMathOperator{\ima}{Im}
\DeclareMathOperator{\id}{Id}
\DeclareMathOperator{\sign}{sign}
\DeclareMathOperator{\ext}{ext}
\newcommand{\ind}{\mathbbm{1}}
\newcommand*{\spr}[2]{\left(\,#1\,\middle|\,#2\,\right)}
\newcommand{\interior}[1]{%
  {\kern0pt#1}^{\mathrm{o}}%
}
\DeclareMathOperator{\card}{card}
\DeclareMathOperator{\ppcm}{ppcm}

%Commandes de pure flemme
\newcommand{\veps}{\varepsilon}
\newcommand{\intab}{\int_{a}^{b}}
\newcommand{\intR}{\int_{-\infty}^{\infty}}
\newcommand{\intinf}{\int_{- \infty}^{+ \infty}}
\newcommand{\equi}{\Leftrightarrow}
\newcommand{\Kev}{$\K$-espace vectoriel }
\newcommand{\Kevs}{$\K$-espaces vectoriels }
\newcommand{\Rev}{$\R$-espace vectoriel }
\newcommand{\Revs}{$\R$-espaces vectoriels }
\newcommand{\Cev}{$\Cx$-espace vectoriel }
\newcommand{\Cevs}{$\Cx$-espaces vectoriels }
\newcommand{\evn}{(E, \norm{.})}
\newcommand{\interf}[2]{[#1,#2]}
\newcommand{\limb}{\lim\limits_}
\newcommand{\lebn}{\lambda_n}
\newcommand{\pp}{\text{~presque partout}}
\newcommand{\derivp}[2]{\frac{\partial #1}{\partial #2}}
\newcommand{\famiu}{(u_i)_{i \in I}}
\newcommand{\sev}{\text{~s.e.v.}}
\newcommand{\Mnp}{M_{n,p}(\K)}
\newcommand{\Mn}{M_{n}(\K)}
\newcommand{\tq}{\text{~t.q.~}}
\newcommand{\mq}{~montrer~que~}
\newcommand{\et}{\quad \text{et} \quad}
\newcommand{\ou}{\text{~ou~}}
\newcommand{\Kx}{\K[X]}


%Commandes de gars chelou
\renewcommand{\subset}{\subseteq}

%Commande restriction, ty duckduckgo
\def\restr#1#2{\mathchoice
              {\setbox1\hbox{${\displaystyle #1}_{\scriptstyle #2}$}
              \restraux{#1}{#2}}
              {\setbox1\hbox{${\textstyle #1}_{\scriptstyle #2}$}
              \restraux{#1}{#2}}
              {\setbox1\hbox{${\scriptstyle #1}_{\scriptscriptstyle #2}$}
              \restraux{#1}{#2}}
              {\setbox1\hbox{${\scriptscriptstyle #1}_{\scriptscriptstyle #2}$}
              \restraux{#1}{#2}}}
\def\restraux#1#2{{#1\,\smash{\vrule height .8\ht1 depth .85\dp1}}_{\,#2}} 

%Quelques normes
\DeclarePairedDelimiter\abs{\lvert}{\rvert}
\DeclarePairedDelimiter\labs{\bigl\lvert}{\bigr\rvert}
\DeclarePairedDelimiter\norm{\lVert}{\rVert}
\DeclarePairedDelimiterXPP\normi[1]{}\lVert\rVert{_1}{\ifblank{#1}{\:\cdot\:}{#1}}
\DeclarePairedDelimiterXPP\normd[1]{}\lVert\rVert{_2}{\ifblank{#1}{\:\cdot\:}{#1}}
\DeclarePairedDelimiterXPP\normp[1]{}\lVert\rVert{_p}{\ifblank{#1}{\:\cdot\:}{#1}}
\DeclarePairedDelimiterXPP\normq[1]{}\lVert\rVert{_q}{\ifblank{#1}{\:\cdot\:}{#1}}
\DeclarePairedDelimiterXPP\norminf[1]{}\lVert\rVert{_\infty}{\ifblank{#1}{\:\cdot\:}{#1}}

%Bracket en angle
\DeclarePairedDelimiter\angl{\langle}{\rangle}

%Set, parenthèses
\newcommand{\set}[1]{\{ #1 \}}
\newcommand{\bigset}[1]{\bigl\{ #1 \bigr\}}
\newcommand{\Bigset}[1]{\Bigl\{ #1 \Bigr\}}
\newcommand{\bigpar}[1]{\bigl( #1 \bigr)}
\newcommand{\Bigpar}[1]{\Bigl( #1 \Bigr)}

%Opitmisation des opérateurs de comparaison
\DeclarePairedDelimiter\tio{o (}{)}
\DeclarePairedDelimiter\gro{O (}{)}


\newcommand{\gast}{\gamma^{\ast}}
\DeclareMathOperator{\Ind}{Ind}

\pagestyle{fancy}

\fancyhead[C]{Analycité des fonctions holomorphes}
\fancyhead[R]{}

\begin{document}

\underline{Source :} Rudin - Analyse - page 244 puis 247-251

\underline{Leçons :}
\begin{itemize}
\item
    230 : Séries de nombres réels et complexes. Comportement des restes ou des sommes partielles des séries numériques. Exemples.
\item
    241 : Suites et séries de fonctions. Exemples et contre-exemples.
\item
    245 : Fonctions holomorphes et méromorphes sur un ouvert de $\Cx$. Exemples et applications.

\end{itemize}

\section{Prérequis topologiques}

\begin{lem}
Si $f$ est une fonction continue de $E$ dans $F$, alors pour tout connexe $C$ de $E$, l'image $f(C)$ est un connexe de $F$.
\end{lem}

\begin{proof}
Posons $C'=f(C)$ et notons $O_1, O_2$ deux ouverts de $F$ tels que $C' \subset O_1 \cap O_2$ 
et $C' \cap O_1 \cap O_2 = \emptyset$. 
On a alors $C \subset f^{-1}(C') \subset f^{-1}(O_1) \cup f^{-1}(O_2)$ avec $ f^{-1}(O_1)$ et $f^{-1}(O_2)$ des ouverts de $E$ par
continuité de $f$. De plus, $C \cap  f^{-1}(O_1) \cap f^{-1}(O_2) = \emptyset$, ce qui entraîne $C \cap  f^{-1}(O_1) = \emptyset$ ou
$C \cap f^{-1}(O_2) = \emptyset$ par connexité de $C$, et donc que $C' \cap O_1 = \emptyset$ ou $C' \cap O_2 = \emptyset$.
L'ensemble $f(C)$ est bien connexe dans $F$.
\end{proof}

\begin{thm}[des fermés emboîtés]
Soit $E$ un espace métrique complet, et $(F_n)_\N$ une suite de fermés emboîtés dans $E$, i.e. :
\[ E \supseteq F_0 \supseteq F_1 \supseteq \dots \]
et dont le diamètre tend vers 0. Alors l'intersection des $F_n$ est réduite à un point.
\end{thm}

\begin{proof}
Pour $n \in \N$, on fixe un point $x_n \in F_n$. On crée ainsi une suite $(x_n)_\N$ qui est de Cauchy.
En effet, comme le diamètre des fermés $F_n$ tend vers 0, pour tout $\veps > 0$, il existe un rang $N \in \N$ tel que ce diamètre 
est majoré par $\veps$. Alors, pour deux points $(x_p,x_q)$ tels que $p, q \geq N$ on a bien $d(x_p,x_q) \leq \veps$.
Comme $E$ est complet, cette suite est convergente et sa limite appartient à tous les fermés $F_n$ 
(car chaque suite $(x_m)_{m \geq n}$ est à valeurs dans $F_n$ fermé et est convergente).
En notant $l$ cette limite, on a $l \in \bigcap_{\N} F_n$. 
Or, pour tout $n \in \N$ le diamètre de $\bigcap_{\N} F_n$ est majoré par le diamètre de $F_n$ qui tend vers 0.
Ainsi, $\bigcap_{\N} F_n$ est réduit au singleton $\set{l}$.
\end{proof}

\section{Prérequis d'analyse complexe}

\begin{defn}[Holomorphie]
Soit $f$ une fonction complexe définie sur un ouvert $\Omega$ du plan. 
Pour $z_0 \in \Omega$, on dit que $f$ est différentiable en $z_0$ si, pour tout $\veps >0$, il existe $\delta >0$ tel que :
\[ \forall z \in D(z_0, \delta), \quad \abs*{ \dfrac{ f(z)-f(z_0) }{ z-z_0 } -f'(z_0) } < \veps . \]
Autrement dit, le quotient $\dfrac{ f(z)-f(z_0) }{ z-z_0 }$ admet une limite en $z_0$ que l'on appelle la dérivée de $f$ en $z_0$ 
(et notée $f'(z_0$)). Si cette dérivée existe pour tout point de $\Omega$, on dit que $f$ est holomorphe sur $\Omega$.
\end{defn}

\begin{lem}
L'ensemble des fonctions holomorphes sur un ouvert $\Omega$ du plan forme un anneau 
(pour l'addition et la multiplication des fonctions).
\end{lem}

\begin{proof}
Les fonctions holomorphes étant en particulier continues, on le déduit en caractérisant séquentiellement les limites sur $\Omega$.
\end{proof}

\begin{thm}
\label{thm:1}
Soit $\mu$ une mesure complexe (finie) sur un espace mesurable $X$, $\varphi$ une fonction complexe mesurable sur $X, \Omega$
un ouvert du plan qui ne rencontre pas $\varphi(X)$, et posons :
\begin{equation}
f(z) = \int_X \dfrac{d\mu(t)}{\varphi(t)-z} \qquad (z \in \Omega) . \label{eq:1}
\end{equation}
Alors, la fonction $f$ est développable en série entière sur $\Omega$.
\end{thm}

\begin{proof}
Prenons un disque $D(a, r) \subset \Omega$. Pour tout $z \in D(a, r)$ et pour tout $t \in X$, on a :
\[ \abs*{ \dfrac{z-a}{\varphi(t)-a }} \leq \dfrac{\abs{z-a}}{r} < 1 . \]
Car $D(a, r) \subset \Omega$ et $\varphi(X)$ ne rencontre pas $\Omega$, donc $\abs{\varphi(t)-a } \geq r$.
On en déduit que la série géométrique :
\[ \sum_{n \in \N} \dfrac{ (z-a)^n }{ (\varphi(t)-a)^{n+1} } = \dfrac1{\varphi(t)-z} , \]
converge uniformément sur $X$, en effet :
\[ \sum_{n \geq N+1}  \dfrac{ (z-a)^n }{ (\varphi(t)-a)^{n+1} } = \left ( \dfrac{ z-a }{\varphi(t)-a} \right )^{N+1} \dfrac1{\varphi(t)-z} , \]
et $\abs*{ \dfrac{z-a}{\varphi(t)-a }}^{N+1} \to_{n \to \infty} 0$. 
On peut donc substituer cette série dans l'intégrale \ref{eq:1} puis intervertir les signes sommes et intégrales, on obtient :
\[  f(z) = \sum_{n \in \N} c_n (z-a)^n \qquad (z \in D(a,r)) , \]
en posant :
\[ c_n = \int_X \sum_{n \in \N} \dfrac{ (z-a)^n }{ (\varphi(t)-a)^{n+1} } \qquad (n \in \N) . \]
\end{proof}

\section{Indices et chemins}

\begin{defn}[Chemin]
On appelle chemin une courbe dans le plan continuement différentiable par morceaux.
Plus précisément, un chemin dont le segment de paramétrage est le segment compact $[a, b] \in \R$ est une fonction $\gamma$,
continue et à valeurs dans $\Cx$ et vérifiant :
\begin{enumerate}[label=(\roman*)]
\item il existe un nombre fini de points $s_j$ tels que $a = s_0 < \dots < s_n = b$, et $\gamma$ a une dérivée continue sur chaque
$[s_j, s_{j+1} ]$ ;
\item aux points $s_1, \dots, s_{n-1}$ (si existants) les dérivées à droite et à gauche de $\gamma$ peuvent différer.
\end{enumerate}
Le chemin est dit fermé si $\gamma(a) = \gamma(b)$.
\end{defn}

\begin{nota}
On nota $\gast$ l'image du chemin $\gamma$.
\end{nota}

\begin{thm}[Indice]
\label{thm:ind}
Soit $\gamma$ un chemin fermé, $\Omega$ le complémentaire de $\gast$ (relativement au plan complexe), et définissons :
\begin{equation}
\Ind_{\gamma}(z) = \dfrac1{2i\pi} \int_{\gamma} \dfrac{dt}{t-z} \qquad (z \in \Omega) . \label{eq:indice}
\end{equation}
La fonction $\Ind_{\gamma}$ est à valeurs entières sur $\Omega$, constante sur chaque composante connexe de $\Omega$ et nulle
sur la composante connexe non bornée de $\Omega$.
\end{thm}

\begin{proof}
Notons $[a,b]$ le paramétrage du chemin $\gamma$ et fixons $z \in \Omega$. Par changement de variable, on a :
\[ \Ind_{\gamma}(z) = \dfrac1{2i\pi} \intab \dfrac{\gamma'(s)}{\gamma(s)-z}ds . \]
Or $w / 2i\pi$ est un entier si, et seulement si, $e^w = 1$. 
Autrement dit, le première assertion du théorème est équivalente à $\varphi(b)=1$, en posant :
\[ \varphi(t) = \exp \left ( \int_a^t \dfrac{\gamma'(s)}{\gamma(s)-z}ds \right ) \qquad (t \in [a,b]) . \]
On peut différentier $\varphi$, sur $[a,b]$ (excepté en les éventuels points $s_1, \dots, s_{n-1}$ définis plus hauts) et on trouve :
\[ \dfrac{\varphi'(t)}{\varphi(t)} = \dfrac{\gamma'(t)}{\gamma(t)-z} . \]
Or $\gamma$ est un chemin fermé, soit $\varphi(b) = 1$ ce qui montre que $\Ind_{\gamma}$ est à valeurs entières sur $\Omega$. \\

D'après le théorème $\ref{thm:1}$, la fonction $\Ind_{\gamma}$ est dévelopable en série entière sur tout $\Omega$, donc
en particulier continue sur cet ensemble. Ainsi, l'image de toute partie connexe de $\Omega$ par $\Ind_{\gamma}$ sera un connexe
de $\N$ soit un singleton ce qui démontre que $\Ind_{\gamma}$ est constante sur chaque composante connexe de $\Omega$.
Enfin, pour $z$ de module suffisament grand, l'expression \ref{eq:indice} est de module strictement inférieur à 1, 
ce qui démontre que $\Ind_{\gamma}(z) = 0$ sur la composante non bornée de $\Omega$.
\end{proof}

\begin{rmq}
L'image $\gast$ étant un compact, il est contenu dans un disque borné $D$ dont le complémentaire $\overline{D}$ est connexe.
Donc $\overline{D}$ est inclus dans une des composantes de $\Omega$, ce qui justifie l'appelation \og la composante connexe
non bornée \fg.
\end{rmq}

\begin{thm}
\label{ind:cercle}
Si $\gamma$ est le cercle orienté positivement de centre $a$ et de rayon $r$, alors :
\[ \Ind_{\gamma}(z) = \left \{ \begin{array}{c c} 1 & \text{  si  } \abs{z-a}<r \\ 0 & \text{  si  } \abs{z-a}>r . \end{array} \right. \]
\end{thm}

\begin{proof}
$\gamma$ est paramétrise sur $[0, 2\pi]$ par $\gamma(t) = a + re^{it}$. 
D'après le théorème $\ref{thm:ind}$, la fonction $\Ind_{\gamma}$ est nulle sur la composante connexe non bornée de $\Omega$
qui est ici le complémentaire du disque. 
Toujours d'après ce théorème, $\Ind_{\gamma}$ est constante sur le disque, on l'estime donc en son centre soit :
\begin{align*}
\Ind_{\gamma}(a) & = \dfrac1{2i\pi} \int_{\gamma} \dfrac{d\zeta}{\zeta-a} \\
& = \dfrac1{2i\pi} \int_0^{2\pi} \dfrac{\gamma'(t)}{\gamma(t)-a} dt \\
& = \dfrac1{2i\pi} \int_0^{2\pi} it dt = 1
\end{align*}
\end{proof}

\section{Théorème local de Cauchy et conséquence}

\begin{thm}
\label{thm:deriv}
Soient $F$ une fonction holomorphe sur un ouvert $\Omega$ et $F'$ une fonction continue sur $\Omega$, 
pour tout chemin fermé $\gamma$ dans $\Omega$, on a :
\[ \int_{\gamma} F'(z) dz = 0 . \]
\end{thm}

\begin{proof}
Notons $[a,b]$ le segment de paramétrage de $\gamma$, comme $\gamma(a) = \gamma(b)$, on déduit :
\[ \int_{\gamma} F'(z) dz = \intab F'(\gamma(t)) \gamma'(t) dt = F(\gamma(b)) - F(\gamma(a)) = 0 . \]
\end{proof}

\begin{coro}
\label{coro:1}
Pour $n \in \Ne$ et pour tout chemin fermé $\gamma$, on a :
\[ \int_{\gamma} z^n dz = 0 . \]
\end{coro}

\begin{proof}
Il suffit de considérer la fonction $z \mapsto z^n$ comme dérivée de la fonction $z \mapsto z^{n+1}/(n+1)$.
\end{proof}

\begin{nota}
Soit $(a,b,c)$ un triplet de nombres complexes, on note $\Delta$ le triangle de sommets $a,b,c$ et pour toute fonction $f$
continue sur la frontière de $\Delta$, on définit l'intégrale :
\[ \int_{\partial \Delta} f = \int_{[a,b]} f + \int_{[b,c]} f + \int_{[c,a]} f . \]
En fait, $\partial \Delta$ est le chemin obtenu en joignant les 3 intervalles.
\end{nota}

\begin{thm}[Cauchy pour un triangle]
\label{cauchy:tgl}
Soit $\Delta$ un triangle fermé, inclus dans un ouvert $\Omega$ du plan.
Soit $p \in \Omega$ et soit $f$ une fonction continue sur $\Omega$ et holomorphe sur $\Omega$ privé de $p$.
On a:
\[ \int_{\partial \Delta} f(z) dz = 0 . \]
\end{thm}

\begin{proof}
\begin{enumerate}[label=(\roman*)]

\item
Supposons d'abord que $p \not \in \Delta$. On note $a', b', c'$ les milieux des segments $[b,c], [c,a]$ et $[a,b]$ respectivement.
On note $\Delta_j$ les quatre triangles formés par les triplets :
\[ \set{a,c',b'}, \quad \set{b,a',c'}, \quad \set{c, a', b'}, \quad \set{a',b',c'} . \]
On a alors :
\[ J = \int_{\partial \Delta} f(z) d = \sum_{j=1}^4 \int_{\partial \Delta_j} f(z) d . \]

\begin{center}
\begin{tikzpicture}

% --- Points ---
\coordinate (a) at (0,0);
\coordinate (b) at (8,0);
\coordinate (c) at (3,4);
\coordinate (a') at (5.5,2);
\coordinate (b') at (1.5,2);
\coordinate (c') at (4,0);

% --- Segments ---
\draw[black, thick] (a) -- (b);
\draw[black, thick] (b) -- (c);
\draw[black, thick] (c) -- (a);

\draw[black] (a') -- (b');
\draw[black] (b') -- (c');
\draw[black] (c') -- (a');

% --- Marques en croix ---
\foreach \P in {a,b,c,a',b',c'}
{
    \draw[black] 
        ($(\P)+(-0.15,-0.15)$) -- ($(\P)+(0.15,0.15)$)
        ($(\P)+(-0.15,0.15)$) -- ($(\P)+(0.15,-0.15)$);
}

% --- Labels ---
\node[above] at (a) {a};
\node[above] at (b) {b};
\node[above] at (c) {c};
\node[right] at (a') {a'};
\node[left] at (b') {b'};
\node[left] at (c') {c'};

\end{tikzpicture}
\end{center} 

Par inégalité triangulaire, la valeur absolue d'au moins une des intégrales du terme de droite est supérieure ou égale à $\abs{J/4}$.
Quitte à réorganiser les indices, supposons que ce triangle est $\Delta_1$. 
On itère alors la procédure ci-dessus sur ce triangle et on forme ainsi une suite $(\Delta_n)_\N$ de triangles tels que :
\[ \Delta \supseteq \Delta_1 \supseteq \Delta_2 \supseteq \dots \]
La longueur du chemin $\partial \Delta_n$ est donnée par $2^{-n}L$, où $L$ désigne la longueur de $\partial \Delta$. 
En outre :
\begin{equation}
\abs{J} \leq 4^n \abs*{ \int_{\partial \Delta_n} f(z) d } \qquad (n \in \Ne). \label{ineq:1}
\end{equation}
D'après le théorème des fermés emboîtés, il existe un unique point commun aux triangles $\Delta_n$ que l'on note $z_0$.
Ce point est aussi un élément de $\Delta$, et $f$ est différentiable en $z_0$. \\
Soit $\veps > 0$, il existe $r > 0$ tel que, pour $z \in D(z_0, r)$ :
\[ \abs*{ \dfrac{ f(z)-f(z_0) }{ z-z_0 } - f'(z_0) } < \veps , \]
ou, de façon équivalente :
\begin{equation}
\abs*{f(z)-f(z_0) - f'(z_0)(z-z_0) } < \veps \abs{z-z_0} . \label{ineq:2}
\end{equation}
De plus, on peut trouver $n \in \N$ tel que $z \in D(z_0, r)$ pour tout $z \in \Delta_n$.
Pour cet entier $n$ on a aussi $ \abs{z-z_0} \leq 2^{-n}L$ pour tout $z \in \Delta_n$.
Or, selon le corollaire $\ref{coro:1}$ on a :
\[ \int_{\partial \Delta_n} f(z_0) + f'(z_0)(z-z_0) dz = 0 , \]
(car linéaire en $z$), on peut donc écrire :
\[  \int_{\partial \Delta_n} f(z) dz =  \int_{\partial \Delta_n} f(z) - f(z_0) - f'(z_0)(z-z_0) dz , \]
de sorte que \ref{ineq:1} entraîne :
\[ \abs*{  \int_{\partial \Delta_n} f(z) dz } \leq \veps (2^{-n}L)^2 . \]
En repartant de l'inégalité \ref{ineq:2}, on obtient :
\[ \abs{J} \leq 4^n \veps 4^{-n} L^2 = \veps L^2 , \]
ceci étant vrai pour tout $\veps >0$, on conclut que $J=0$ pour $p \not \in \Delta$.

\item
Supposons maintenant que $p$ est un sommet de $\Delta$, par exemple $p=a$. 
Si les points $a, b, c$ sont alignés, alors le résultat est trivialement nul (on intègre en faisant un \og aller-retour \fg{} sur un segment). \\
Sinon, on choisit $x \in [a,b]$ et $y \in [a,c]$ tous les deux voisins de $a$ et on observe que l'intégrale de $f$ sur $\partial \Delta$
est la somme des intégrales sur les frontières des trois triangles $\set{a,x,y}, \set{b,x,y}$ et $\set{b,c,y}$.
\begin{center}
\begin{tikzpicture}

% --- Points ---
\coordinate (a) at (0,0);
\coordinate (b) at (8,0);
\coordinate (c) at (3,4);
\coordinate (x) at (2,0);
\coordinate (y) at (0.6,0.8);

% --- Segments ---
\draw[black, thick] (a) -- (b);
\draw[black, thick] (b) -- (c);
\draw[black, thick] (c) -- (a);

\draw[black] (x) -- (y);
\draw[black] (b) -- (y);

% --- Marques en croix ---
\foreach \P in {a,b,c,x,y}
{
    \draw[black] 
        ($(\P)+(-0.15,-0.15)$) -- ($(\P)+(0.15,0.15)$)
        ($(\P)+(-0.15,0.15)$) -- ($(\P)+(0.15,-0.15)$);
}

% --- Labels ---
\node[left] at (a) {a};
\node[above] at (b) {b};
\node[above] at (c) {c};
\node[below] at (x) {x};
\node[left] at (y) {y};

\end{tikzpicture}
\end{center} 
Ces deux dernières intégrales sont nulles : les triangles ne contenant pas $p$, on leur applique le résultat (i).
Il en résulte que :
\[ \int_{\partial \Delta} f(z) d = \int_{[a,x]} f(z) dz + \int_{[x,y]} f(z) dz + \int_{[y,a]} f(z) dz , \]
Chacun de ces intervalles peut être rendu arbitrairement petit et $f$ est bornée sur chacun de ces intervalles (car continue), 
donc l'intégrale de $f$ sur le chemin $\partial \Delta$ est nulle.

\item
Finalement, supposons que $p$ est un point quelconque du triangle (hors sommet), alors on applique le point (ii) aux trois triangles
$\set{a,b,p}, \set{b,c,p}$ et $\set{a,c,p}$ ce qui achève la démonstration.

\end{enumerate}
\end{proof}

\begin{thm}[Cauchy pour un ensemble convexe]
\label{Cauchy:conv}
Soit $\Omega$ un ouvert convexe, $p \in \Omega$, $f$ une fonction continue sur $\Omega$ et holomorphe sur $\Omega$
privé de $p$. Pour tout chemin fermé $\gamma$ dans $\Omega$ :
\[ \int_{\gamma} f(z) dz = 0 . \]
\end{thm}

\begin{proof}
On fixe $a \in \Omega$. Pour tout $z\in \Omega$, on a $[a,z] \in \Omega$ par convexité, de sorte qu'on peut définir :
\[ \forall z \in \Omega, \quad F(z) ) \int_{[a,z]} f(t) dt . \]
Pour tous $z, z_0 \in \Omega$, le triangle $\Delta = \set{a, z, z_0}$ es inclus dans $\Omega$, donc :
\[ \int_{\partial \Delta} f(t) dt = 0 = \int_{[a,z]} f(t) dt + \int_{[z,z_0]} f(t) dt + \int_{[z_0,a]} f(t) dt , \]
d'après le théorème $\ref{cauchy:tgl}$. On en déduit que :
\[ \int_{[z_0, z]} f(t) dt = F(z) - F(z_0) . \]
En fixant $z_0 \in \Omega$ et en considérant $z \neq z_0$, on obtient :
\[ \dfrac{F(z)-F(z_0)}{z-z_0} - f(z_0) = \dfrac1{z-z_0} \int_{[z_0, z]} (f(t) - f(z_0)) dt . \]
Soit $\veps >0$, par continuité de $f$ en $z_0$, il existe $\delta >0$ tel que :
\[ \forall t \tq \abs{t-z_0} < \delta, \quad \abs{f(t)-f(z_0)} < \veps ; \]
autrement dit, $F'(z_0) = f(z_0)$ pour tout $z_0 \in \Omega$ et on déduit le résultat du théorème $\ref{thm:deriv}$.
\end{proof}

\begin{thm}[formule de Cauchy pour un ensemble convexe]
\label{form:Cauchy}
Soit $\gamma$ un chemin fermé dans un ouvert convexe $\Omega$, et soit $f$ une fonction holomoprhe sur $\Omega$. 
Si $z \in \Omega$ et si $z \not \in \gast$, on a :
\[ f(z) \Ind_{\gamma}(z) = \dfrac1{2\pi i} \int_{\gamma} \dfrac{f(t)}{t-z} dt . \]
\end{thm}

\begin{proof}
Soit $z$ tel qu'énoncé et définissons :
\[ g(t) = \left \{ \begin{array}{c c}
\dfrac{ft)-f(z)}{t-z} & \text{  si  } t \in \Omega, t \neq z \\
f'(z) & \text{  si  } t = z .
\end{array} \right . \]
La fonction $g$ satisfait aux hypothèses du théorème $\ref{Cauchy:conv}$ d'où :
\[ \dfrac1{2\pi i} \int_{\gamma} g(t) dt = 0 , \]
autrement dit, comme $z \not \in \gast$ :
\begin{align*}
\dfrac1{2\pi i} \int_{\gamma} \dfrac{ft)-f(z)}{t-z} dt & = 0 \\
& = \dfrac1{2\pi i} \int_{\gamma} \dfrac{f(t)}{t-z} dt - \dfrac1{2\pi i} \int_{\gamma} \dfrac{f(z)}{t-z} dt \\
& = \dfrac1{2\pi i} \int_{\gamma} \dfrac{f(t)}{t-z} dt - f(z) \Ind_{\gamma}(z)
\end{align*}
\end{proof}

\begin{thm}
Pour tout ouvert $\Omega$ du plan, toute fonction holomorphe sur $\Omega$ est développable en série entière dans $\Omega$.
\end{thm}

\begin{proof}
Soient $f$ une telle fonction et $D(a, R) \subset \Omega$. 
Si $\gamma$ est un cercle orienté positivement, de centre $a$ et de rayon $r < R$, la convexité de $D(a, R)$ nous permet d'appliquer
la formule de Cauchy (cf. $\ref{form:Cauchy}$) soit :
\[ f(z) \Ind_{\gamma}(z)  = \dfrac1{2\pi i} \int_{\gamma} \dfrac{f(t)}{t-z} dt  \qquad (z \in D(a,r)) \]
et $\Ind_{\gamma}(z) = 1$ pour $z \in D(a,r)$ selon le théorème $\ref{ind:cercle}$.
Par changement de variable, on obtient :
\[ f(z) = \int_0^{2\pi} \dfrac { f(\gamma(t)) \gamma'(t) dt }{\gamma(t)-z} \qquad (z \in D(a,r)) \]
On applique maintenant le théorème $\ref{thm:1}$ en prenant $X= [0, 2\pi], d\mu(t) = f(\gamma(t)) \gamma'(t) dt$ 
et $\varphi = \gamma$. 
Les hypothèses du théorème sont satisfaites et on conclut que $f$ est développable en série entière sur $D(a,r), r < R$.
L'unicité d'un tel développement, assure que l'on obtient la même représentation pour tout $r < R$, 
donc que ce développement en séries entières est valable pour tout $z \in D(a, R)$, ce qui termine la démonstration.
\end{proof}

\begin{rmq}
La fonction $f$ étant holomorphe sur $\Omega$ (en particulier continue), en définissant :
\[ F(w) = \int_a^w f(t) dt \qquad (w \in \Omega) . \]
La fonction $F \circ \gamma$ définie bien une mesure complexe sur $[0, 2\pi]$.
\end{rmq}


\end{document}